\section{Discussion}

\subsection{Positioning of the Stage-3 Surrogate as a Screening Model}

The synthetic dataset is generated by the reduced-order \(\mathrm{SGQR}\) fitting in Section~\ref{sec:SGQR-target-construction-en}: raw CDF trajectories are first subjected to cleansing and temporal alignment, then fitted to the \(\mathrm{SGQR}\) form, providing smooth extrapolation targets within the specified time window.

A natural objection at this point is why the workflow does not stop here and simply promote the SGQR fit itself (or an improved empirical correlation built from its scaling exponents) to a stand-alone predictive model, rather than proceeding through three additional training stages to outperform Hiroyasu--Arai directly. Chapter~4's own diagnostics of the SGQR fit answer this. High per-trajectory fit accuracy (median RMSE \(\SI{1.27}{mm}\) on accepted fits) only certifies that the three-parameter sigmoid-gated quarter-root form can retrace an already-observed trajectory; it does not certify that \((k_{1/4},t_0,s)\) are transferable, physically meaningful functions of operating condition. The \(90.0\%\) post-fit-shielding pass rate is, by construction, a fit-quality and shape-plausibility gate (a bounded far-time plausibility window, parameter sanity checks, and robust per-file z-scores on RMSE and cost-per-point), not a test of physical generalization: a trajectory passes because the monotonic sigmoid-gated shape happens to trace it well, independent of whether the underlying breakup and evaporation physics is correctly represented. Most fundamentally, the identifiability analysis (Figure~\ref{fig:SGQR-identifiability}) shows that \(k_{1/4}\), \(t_0\), and \(s\) are not independent physical knobs: median pairwise correlations of \(\sim0.76\) and \(\sim0.56\) indicate a softly redundant, coupled three-dimensional surface rather than a clean decomposition into amplitude, onset time, and gating width, and the pooled scaling regression built on top of this coupled surface is itself weak (\(R^2\approx0.27\)) and confounded by nozzle-diameter/pressure collinearity in the design matrix (Section~\ref{sec:pointlevel-scaling-verification}). Taken together, these findings mean SGQR is fit for its intended purpose: compressing noisy plume trajectories into smooth, differentiable, per-trajectory targets, but its fitted parameters carry no validated claim to physical generalization, so the fit cannot itself be published as a replacement empirical law for Hiroyasu--Arai. Predictive credibility is instead earned downstream, by training a surrogate on the full population of these fits and validating it end-to-end against raw CDF observations, which is precisely the role assigned to Stages~1--3 below.

Stage~1 extracts stable representative mean trajectories condition by condition from the complete synthetic dataset, validating the $A$-scaling hypothesis.
However, because only one representative trajectory is retained per operating condition, using Stage~1 in isolation neglects plume-to-plume and shot-to-shot variability, rendering it unsuitable as a definitive screening model.

Stage~2 extends the supervision to the entire synthetic dataset, concurrently learning the conditional mean and conditional dispersion within the extrapolation window via two output heads.
This also explains its slightly higher MAE compared to Stage~1: Stage~2 faces a harder supervision problem that more closely matches the final screening requirement.

Stage~3 is designed around selecting different supervision signals based on the temporal period.
In early periods, where measured data are relatively dense, the model learns directly from observed CDF data, preserving the details of the true distribution.
In late periods, where right-censoring makes observations sparse, the model instead aligns with the conditional mean and variance from the Stage~2 teacher model, following the rationale established in Chapter~2.
The two supervision signals are combined through a mixed KD loss that consists of a mean MSE term and a log-variance MSE term (see Section~\ref{sec:kd-form-fix-en}).
Because both the teacher and student output unimodal Gaussian distributions, this loss has a closed-form solution, so Monte Carlo sampling is not needed and no additional sampling noise is introduced.
This hybrid supervision explains why Stage~3 shows higher fidelity in early periods and more conservative behavior in late periods.

\subsection{Synthesis of Validated Results and Boundaries of Applicability}

The validation evidence from Chapters~4 and~5 supports four principal interpretations.
First, finite-window censoring is a substantial property of the archive: 39.6\% of the cleaned CDF penetration trajectories reach their nozzle-specific field-of-view limit.
The three-parameter \(\mathrm{SGQR}\) fit addresses the resulting target-construction problem more robustly than the four-parameter two-branch candidate, increasing the post-fit-shielding pass rate from 61.3\% to 90.0\% while retaining a median accepted-fit RMSE of \SI{1.27}{mm}.
This evidence supports its use as a smooth trajectory compressor and target generator, but does not validate physical continuation beyond the camera window.

Second, the time-resolved scaling analysis shows that the effective pressure-difference exponent is close to 0.5 over the well-supported early-to-mid transient window before moving towards the classical post-breakup asymptote.
Relative to the \(\Delta P^{0.25}\) alternative, the adopted \(\Delta P^{0.5}\) scaling produces an approximately \(6.6\times\) tighter collapse of the representative trajectory population.
This result provides the empirical basis for the amplitude scaling and residual-pressure feature design used by the surrogate; its physical interpretation is developed further in Section~\ref{sec:time-varying-scaling-discussion-en}.

Third, on the primary uncensored CDF protocol (540{,}367 held-out points per seed from the nine-nozzle archive), the production MLP achieves a five-seed mean RMSE of \SI{4.185}{mm} and MAE of \SI{2.978}{mm}.
Under this dataset, protocol, and the refitted baseline definitions, this is a 58.9\% RMSE reduction relative to Hiroyasu--Arai (\SI{10.192}{mm}) and a 54.6\% reduction relative to Naber--Siebers (\SI{9.212}{mm}).
On the same fixed table, the seed-42 SVGP reference attains lower RMSE and CRPS (\SI{4.087}{mm} and \SI{2.063}{mm}, versus \SI{4.185}{mm} and \SI{2.301}{mm} for the MLP), whereas the MLP produces wider in-domain intervals (\(1\sigma/2\sigma\) coverage 0.888/0.990 versus 0.717/0.946) and supports the onset output and regime-aware distillation chain used by the screening workflow.
Because the seed protocols are unmatched, this comparison is descriptive rather than a statistical-significance claim, and conservative coverage is not a false-negative safety guarantee.

Fourth, the five-fold LONO study defines a clear transfer boundary.
After excluding Nozzle~0, the four-fold mean RMSE is \SI{7.00}{mm} for the MLP and \SI{5.95}{mm} for the SVGP, indicating useful but imperfect transfer within the evaluated design lineage.
On Nozzle~0, the corresponding errors rise to approximately \SI{35}{mm} and \SI{27}{mm}.
Nozzle~0 is a factory-fresh reference injector outside the prototype design lineage, so this fold combines operating-condition extrapolation with an injector-population shift; it should not be interpreted as an ordinary in-family fold or generalized into a claim about all unseen hardware families.

Table~\ref{tab:discussion-validated-results} consolidates the quantitative evidence and the boundary attached to each interpretation.

\begin{table}[t]
\centering
\footnotesize
\setlength{\tabcolsep}{4pt}
\begin{tabularx}{\textwidth}{L{3.2cm} Y L{4.2cm}}
\toprule
Finding / Theme & Quantitative Evidence & Operational Scope / Boundary \\
\midrule
FOV Right-Censoring & 39.6\% of cleaned CDF trajectories reach the FOV limit & Common to the finite-window Mie archive \\
SGQR Fit Robustness & Pass rate improved from 61.3\% to 90.0\% (median RMSE \SI{1.27}{mm}) & Trajectory compressor; wide-FOV continuation unvalidated \\
Transient Pressure Scaling & \(\Delta P^{0.5}\) gives \(\sim6.6\times\) tighter collapse than \(\Delta P^{0.25}\) & Supported early-to-mid transient window \\
Empirical Baseline Gain & MLP RMSE reduced by 58.9\% (H--A) and 54.6\% (N--S) & Primary protocol and refitted baselines only \\
Model Trade-off & SVGP RMSE/CRPS: \SI{4.087}{mm}/\SI{2.063}{mm}; MLP: \SI{4.185}{mm}/\SI{2.301}{mm} with wider coverage & Descriptive comparison with unmatched seed protocols \\
LONO Transfer Boundary & Excluding Nozzle~0: MLP \SI{7.00}{mm}, SVGP \SI{5.95}{mm}; Nozzle~0: MLP approx.\ \SI{35}{mm}, SVGP approx.\ \SI{27}{mm} & Nozzle~0 combines extrapolation and injector-population shift \\
\bottomrule
\end{tabularx}
\caption{Synthesis of the primary validated results, their quantitative evidence, and their operational boundaries.}
\label{tab:discussion-validated-results}
\end{table}

The intended industrial value of this workflow is as an early screening layer that reuses existing optical archives at low cost, narrows the set of nozzle and operating-condition candidates advanced to CFD simulation or bench testing, and flags geometrically high-risk cases for closer inspection, though the magnitude of that narrowing has not itself been measured in this thesis (see the cost-savings caveat below).
Its scope should be stated clearly.
Quantifying those savings is an engineering task for the deployment phase, once the screening layer is integrated into an industrial decision process and operational data have accumulated; it falls outside this thesis's scope.

The extracted spray boundaries represent the liquid-phase envelope observed in Mie imaging under a specific optical protocol, not the complete in-cylinder liquid-phase ground truth \cite{ecn_liquid_penetration,ecn_liquid_penetration_accessory}.
The model is primarily tailored for interpolative prediction within the training operating condition domain, as opposed to unconstrained extrapolation;
moreover, it is not intended to supplant CFD, engine testing, or validated physical models of spray-wall impingement, but rather acts as an auditable early-stage engineering screening tool to narrow the design space before initiating higher-fidelity validation.

The target definition rests on shape priors rather than physical mechanisms: the reduced-order fitting model extrapolates each spray penetration trajectory in time using growth dominated by the quarter-power law ($t^{1/4}$), and the neural surrogate is constrained to be monotonically increasing via regularization; neither is derived from droplet evaporation or breakup physics.
These shape priors make trajectory growth intuitively reasonable and correct short-term extrapolations at the tail end of the observation window, but they do not model the imaging mechanism of Mie scattering.
This thesis further hypothesizes that, in room-temperature, non-reacting Mie experiments, the late-stage attenuation of the image signal primarily originates from droplet breakup, spatial diffusion, localized reduction in droplet number density, scattering intensity dropping below the detection threshold, and detection side-effects induced by the brightness multiplier mask. Rapid evaporative disappearance of the fuel within a millisecond-scale window is not the primary driver.
Consequently, even if the plume leading edge progressively becomes invisible before reaching the boundary of the observation window, this thesis still interprets it as a decline in the visibility of the liquid-phase geometric envelope, rather than the physical termination of droplet motion.
This interpretation accords with the physical intuition of relatively slow diesel evaporation at room temperature, though this thesis lacks synchronous measurements of droplet size, mass flux, or evaporation rate to validate it directly.
Thus, the derived penetration should be construed as the continuation of the Mie-defined geometric envelope, rather than direct evidence of a robust liquid-phase spray tip, localized liquid mass, droplet size, phase state, or the potential for wall film formation \cite{ecn_liquid_penetration,ecn_liquid_penetration_accessory}.

This extrapolation hypothesis still lacks direct comparative validation.
The originally planned validation pathway was to use wide-field single-plume Mie videos with an approximately twofold field of view, identical nozzle types, and similar operating conditions provided by an industrial collaborator, to comprehensively observe the penetration trajectory until the plateau region;
subsequently, the field of view would be artificially truncated, fitting the \(\mathrm{SGQR}\) model solely with the pre-truncation time series, and comparing the extrapolated curve pointwise with the full measured trajectory that remains visible post-truncation.
This experiment would directly scrutinize the core hypothesis of the knowledge distillation (KD) chain in this thesis: that the parametric form \(\sigma\!\left((t-t_0)/s\right)\cdot k_{1/4}\,t^{1/4}\) remains a credible extension beyond the boundary of the field of view, rather than a deviation primarily governed by shape priors.
Due to the collaborator's inability to provide such single-plume data, this validation could not be completed within the scope of this thesis; therefore, the reliable time range for extrapolation using the \(\mathrm{SGQR}\) hypothesis remains an open empirical question.

The aforementioned boundaries are critical for downstream wall-impingement screening.
The interaction between the spray and the wall is not invariably deleterious: if the majority of the liquid phase has evaporated prior to reaching the wall, the residual wall interaction might be less detrimental, and under certain conditions, it might even facilitate entrainment and mixing.
What truly warrants prioritized attention is direct liquid droplet impingement and continuous wall film formation \cite{MorieiraMotaPanao2010,liu2024marine_impingement,peraza2022swi,ahmad2025fueldilution}.
From this perspective, the present room-temperature, non-reacting, non-evaporating Mie dataset can be interpreted as a conservative stress case for liquid-envelope persistence: if a candidate spray remains well separated from representative metal surfaces under these cold optical conditions, faster evaporation in an engine-representative hot environment would generally strengthen, rather than weaken, the early low-risk indication.
This argument is useful for upstream screening, but it is not a proof of engine safety.
Real in-cylinder operation also introduces moving boundaries, pressure waves, squish, swirl, crossflow, temperature gradients, and spray deflection; these mechanisms can change the relative path between the liquid envelope and nearby metal surfaces even while evaporation shortens the liquid lifetime.
Accordingly, the model in this thesis should be more accurately interpreted as a conservative screening surrogate for the probable proximity of the Mie-defined spray envelope to the wall, rather than a phase-resolved predictor of harmful liquid-state impingement.
The correspondence between the liquid-phase envelope geometry and temperature, evaporation, local droplet mass flux, and wall film formation has not been modeled or validated in this thesis;
an elevated wall-interaction score merely denotes a higher risk of geometric proximity, and cannot be directly equated with harmful liquid impingement or the quantity of wall wetting.
Likewise, conservative coverage of penetration uncertainty is not identical to false-negative safety: a ranking tool can be cautious in its geometric envelope while still missing harmful wall wetting if evaporation, droplet inertia, cone-angle evolution, or in-cylinder flow changes the actual liquid mass reaching the surface.

The downstream engine geometric coupling is similarly a simplified model.
The current prototype exclusively translates the piston wall mask using slider-crank kinematics, without back-incorporating time-varying chamber pressure, pressure wave dynamics, or large-scale bulk gas motion into the penetration prediction.
This simplification is inherited from the OSCC measurement itself, as every injection event occurs in a rigid chamber with a fixed back pressure;
consequently, what the surrogate learns is the penetration under quasi-static chamber conditions, rather than the fully coupled in-cylinder transient state.
Considering the sensitivity of diesel spray penetration to ambient gas density, this decoupling should be apprehended as a pragmatic modeling choice, rather than a proof that pressure and gas flow effects are negligible \cite{NaberSiebers1996}.
In the target scenario of this thesis, namely short pilot injection windows near top dead center on a millisecond timescale, this approximation remains acceptable as a first-order engineering screen: under the current toy piston setup, the crown displacement is less than \SI{1}{mm} within a \SI{5}{ms} window, and the present requirement is geometric ranking, not crank-angle-resolved in-cylinder ground truth prediction.
When the engine speed increases, the injection timing deviates further from top dead center, or the in-cylinder transient pressure and gas velocity significantly alter the ambient conditions perceived by the spray, the current wall-interaction score can only be construed as a decoupled screening proxy.

Another pertinent limitation relates to the spray cone angle.
Although a cone-angle proxy was extracted during the image processing phase, this thesis neither systematically analyzed it across different campaigns, nor elevated it to a learning target for the final surrogate.
Therefore, the half-cone angle \(\alpha\) in the downstream wall-interaction model is a fixed input, rather than a validated, time-varying observable.
Existing spray literature indicates that the spreading angle varies with transient development, ambient gas density, injection conditions, and nozzle geometry \cite{NaberSiebers1996,JungManinSkeenPickett2015}.
The current geometric screening solely learns the uncertainty in the axial penetration direction; the transverse envelope determined by \(\alpha\) remains postulated rather than predicted.

The remaining boundaries fall into three categories: generalization and out-of-distribution limits, statistical and sample-independence caveats, and comparative-baseline scope.

\subsubsection*{Generalization and Out-of-Distribution Bounds}
\begin{itemize}
	\item The surrogate is strictly trained within the operating condition domain jointly supported by experimental and synthetic data.
	      Under genuinely out-of-distribution pressures, nozzle geometries, ambient temperatures, or injection strategies, its predictions should not be accorded equivalent credibility.

	\item The leave-one-nozzle-out evaluation is dominated by the Nozzle0 fold (the sole factory-fresh reference nozzle, with injector-to-injector variation not modeled), on which every Stage~3 variant collapses (RMSE \(\sim\SI{35}{mm}\), systemic underestimation; per-fold breakdown in Figure~\ref{fig:svgp-lono-comparison}); the five-fold mean RMSE of \SI{12.55}{mm} is therefore a worst-case figure carried by this one out-of-design-family fold.
	      With that fold excluded, the in-family four-fold mean (\SI{7.00}{mm}) falls below both physical baselines under the identical protocol (Chapter~5). The interpretive point is that the model's cross-family risk is localized to genuinely off-design geometry, not to in-family penetration prediction.
	      The image evidence of Nozzle~3 (Figure~\ref{fig:failure-mode-nozzle3}) further correlates the residual errors within the design family to leading liquid slugs, chasing and overtaking main jets, inter-hole asymmetry, and step-wise penetration. This trajectory morphology is consistent with the wake-driven chasing/overtaking mechanism reported in the split/multiple-injection literature.
	      However, as this thesis possesses neither synchronous gas velocity measurements nor internal flow visualization or liquid slug identity tracking, it can only interpret Nozzle~3 as image evidence consistent with this mechanism; mechanical causes such as needle wobble or transient nozzle hole throttling remain equally viable candidate explanations that the data cannot exclude.
	      The principal cross-family risk of the current model is the inadequate generalization to geometries genuinely outside the design family (Nozzle0), rather than a holistic failure to support average penetration prediction within the design family.
\end{itemize}

\subsubsection*{Statistical and Sample-Independence Constraints}
\begin{itemize}
	\item Seed-to-seed RMSE is stable across the five anchor-off repeated training runs (std \SI{0.039}{mm}, Chapter~5), so the headline mean is not a single-seed artefact.
	      It is still not a full statistical certification, however: the refreshed fixed-table CSV reports cluster-bootstrap intervals for RMSE, MAE, and bias, while condition-level confidence intervals remain to be established.

	\item The held-out RMSE and MAE are computed based on 540,367 point-wise samples, but the temporal point errors within the same trajectory are mutually correlated.
	      Thus, the number of independent samples is significantly smaller than the number of point-wise samples, and any confidence statements based on point-wise counts should be understood in conjunction with this correlation.
	      Condition-level confidence intervals, preferably clustered by recording, condition, and nozzle, are still required before the tool can support formal acceptance thresholds.
\end{itemize}

\subsubsection*{Comparative-Baseline Boundaries}
\begin{itemize}
	\item Two categories of baselines are included in the comparison: physical empirical correlations (H-A and N-S) and a non-parametric probabilistic baseline (SVGP). Under the same evaluation protocol, both machine learning methods substantially outperform the physical correlations. Between the two, SVGP achieves slightly higher point accuracy (RMSE \SI{4.087}{mm} vs.\ \SI{4.185}{mm}), while the MLP produces more conservative uncertainty coverage ($1\sigma$/$2\sigma$ = 88.8\%/99.0\% vs.\ 71.7\%/94.6\%).

	      Lightweight regressors such as random forest, gradient boosting, or NARX have not been systematically included in the comparison; this thesis therefore makes no claim that the MLP's point accuracy is superior to all machine learning alternatives.

	      The choice of MLP, however, is not based on point accuracy. It rests on three structural capabilities: (1)~a joint mean-and-uncertainty output head that directly supports the probabilistic screening interface; (2)~a knowledge distillation chain that mixes synthetic right-censored targets with raw observations during forward propagation; and (3)~a lightweight backbone that makes a full grid of 60+ ablations $\times$ LONO $\times$ multiple training seeds feasible. None of these capabilities can be replicated at low cost by random forests, gradient boosting, or NARX. Without any one of them, the design rationale for the screening interface does not hold.

	\item This chapter's censoring rationale for the three-stage curriculum (Section~\ref{sec:mlp-curriculum-en}) is not confirmed by a controlled ablation that trains the surrogate end-to-end directly on the raw, censored Dataset~\#1 without the SGQR-fitted intermediate stages. The censoring audit of Section~\ref{sec:cdf-spatial-censoring-audit} motivates the expectation that such a single-stage baseline would exhibit saturation bias near the field-of-view boundary, but this expectation is not empirically demonstrated within the present computational budget and is left as future work.
\end{itemize}

\subsubsection*{MLP vs.\ SVGP: Accuracy Trade-off and Model Selection}

The accuracy trade-off is established in Chapter~5 and synthesized in Table~\ref{tab:discussion-validated-results}: on raw point accuracy, the SVGP holds a marginal advantage (RMSE \SI{4.087}{mm} vs.\ \SI{4.185}{mm}, a gap of \SI{0.098}{mm}) and transfers more gracefully to unseen nozzles under LONO, while the MLP leads on uncertainty coverage (1$\sigma$/2$\sigma$: $0.888/0.990$ vs.\ $0.717/0.946$). The MLP was therefore selected not for overall accuracy superiority, but for three structural capabilities the kernel baseline cannot easily provide: the onset-CDF output head required by the screening interface (Section~\ref{sec:screening-interface-en}), the regime-aware censoring chain that mixes fitted synthetic targets with right-censored raw observations, and the design-axis ablations covering stage curriculum and KD form. These structural choices do not exist in the SVGP framework and therefore cannot be ablated there (Section~\ref{sec:mlp-ablation-affordability-en}).

The optional per-family residual adapter illustrates a possible extension of this payoff.
Adding a small correction head to the frozen trunk reduces in-family LONO RMSE by
approximately 31\% ($7.44 \rightarrow \SI{5.16}{mm}$) and removes the mean bias (Chapter~5).
The limit is equally clear: a nozzle outside the training design space cannot be
recovered by a lightweight adapter: Nozzle0 plateaus well above the \SI{5}{mm}
in-family band, marking the point at which new trunk capacity is required rather than
a new head. Kernel baselines do not naturally support the same adapter paradigm: adapting an SVGP to a
new family requires re-optimising the full set of variational parameters and inducing
points, which is closer to a full retrain than a lightweight fine-tune.

\subsection{Empirical Physical Finding: Time-Dependent Pressure Scaling in Short Pilot Injections}
\label{sec:time-varying-scaling-discussion-en}

The time-windowed log-log regression in Chapter~4 (Figure~\ref{fig:time-windowed-exponents}) provides the thesis's main empirical physical finding: within the observable window of short pilot injections, a single global pressure-difference exponent does not describe the full archive. On the Level-3 QC-gated, censor-truncated population, the pooled effective \(\Delta P\) exponent spans approximately 0.45--0.69 over the high-support \SIrange{0.3}{0.7}{ms} interval and then moves towards the Hiroyasu--Arai post-breakup reference value. The trajectory is not strictly monotone---for example, the estimate rises locally between 0.5 and \SI{0.6}{ms}---so it is evidence of an overall time dependence, not a fitted switching law.

The early-to-mid estimates retain a Bernoulli/orifice-velocity imprint, whereas the later pooled values are consistent with movement towards the H--A asymptote. That interpretation is physically plausible, but the analysis does not identify a unique regime-transition time. Under the censor-truncated protocol, \(a_{\Delta P}\) reaches about 0.30 at \SI{0.8}{ms} and 0.23 at \SIrange{0.9}{1.0}{ms}; at the same time, the represented condition count falls from 589 at \SI{0.3}{ms} to 177 at \SI{1.0}{ms}, and point-level \(R^2\) falls to about 0.09. The censoring-inclusive curve approaches the same range later. A fixed cohort of conditions that remains observable at \SI{1.0}{ms} also trends downward, but begins near 0.41 rather than the pooled 0.69. Thus, part of the apparent evolution reflects changing observational support as well as physical time evolution.

The regression also retains a measurable ambient-pressure coefficient conditional on \(\Delta P\) and the diameter control. It is approximately \(-0.23\) to \(-0.25\) over \SIrange{0.3}{0.5}{ms}, which is consistent with the classical ambient-density exponent and shows that a \(\Delta P\)-only scaling is incomplete in this archive. Because physical chamber pressure and ambient density are proportional under the present constant-temperature mapping, the regression cannot separate these two descriptions. Nor can it attribute the partial coefficient causally to cavitation, throttling, or internal nozzle flow: changing \(P_\mathrm{ch}\) while holding \(\Delta P\) also changes \(P_\mathrm{inj}\), and the nozzle/pressure support is confounded.

Together, these results show that any fixed exponent used by the surrogate---whether the H--A \(\Delta P^{0.25}\) term or the adopted \(\Delta P^{0.5}\) term---is a mean-field preconditioner valid only for a specified temporal and conditional domain. The feature-representation LONO ablation supports the predictive utility of \(\Delta P^{0.5}\) and the residual pressure channels within the present data distribution; it does not independently prove the proposed mechanism. The appropriate physical claim is therefore a censoring-aware, time-resolved empirical scaling result, not a new universal penetration law.

This finding does not invalidate the Hiroyasu--Arai relation as a post-breakup asymptotic law. It shows instead that the observable mass of these short-pilot-injection records includes a substantial transitional interval in which the post-breakup asymptote is not a sufficient global description. That distinction is both the physical result and the empirical basis for the surrogate's feature design.

\subsection{Future Validation Pathways for the Probabilistic Wall-Impingement Screening Prototype}
\label{sec:future-validation-pathways-en}

The wall-impingement screening prototype delineated in Section~\ref{sec:screening-interface-en} realizes the probabilistic proxy advocated in the literature review, but the frame-by-frame impingement probability score \(P_{\mathrm{hit}}(t)\) has not yet been benchmarked against physical spray-wall observations.
This gap defines a specific next validation step because the required reference data are already available in the industrial collaborator's experimental dataset.

For industrial deployment, five criteria should be satisfied before the prototype is used as anything stronger than a rank-oriented pre-screen.
First, the \(\mathrm{SGQR}\)-based late continuation should be validated against wide-field single-plume Mie recordings by artificially truncating the field of view and comparing the extrapolated trajectory with the full visible measurement.
Second, engine-representative optical wall data should establish whether the geometric interception score detects the first observed wall-contact frame and ranks cumulative spray exposure correctly.
Third, a CFD-annotated benchmark should quantify screening recall, false-negative rate, and computational savings for the proposed triage workflow.
Fourth, uncertainty summaries should be reported at condition or campaign level with independent confidence intervals rather than relying only on correlated point-wise samples.
Fifth, the fixed half-cone angle should be replaced or bounded by a learned, time-varying cone-angle model once the cone-angle proxy extracted in Chapter~3 has been validated.

The most direct validation target is high-speed optical experiments under representative piston geometries, namely, frame-by-frame imaging of the injection process under engine-representative boundary conditions.
Such experiments can directly observe the expansion of the spray in the vicinity of the piston crown and compare it with \(P_{\mathrm{hit}}(t)\):
The first frame in which spray contact is observed provides an empirical boundary for impingement onset, and the cumulative spray exposure within the injection window can be compared with the integrated score \(E = \int P_{\mathrm{hit}}\,\mathrm{d}t\).
This comparison will scrutinize the axial penetration surrogate under dynamic geometric conditions that the training dataset cannot provide.

Another complementary pathway is CFD numerical simulation.
High-fidelity spray calculations can resolve liquid-phase distribution, wall liquid film mass, and local droplet concentration, whereas the surrogate defined by Mie imaging merely approximates these physical quantities with a rotated Gaussian envelope.
The positioning of this surrogate is not to supplant CFD, but to serve as a low-cost upstream screening layer: operating conditions where the impingement probability exceeds the threshold proceed to full computation, while those safely below the threshold are filtered out.
Validating the screening recall rate, false negative rate, and computational savings of this stratified workflow necessitates a set of CFD-annotated reference operating conditions; the current dataset does not yet possess these.
Only after such a benchmark can conservative geometric coverage be translated into a defensible industrial safety margin.

Both validation pathways will resolve a structural open question in the current prototype: whether a fixed half-cone angle \(\alpha\) is acceptable, or whether it is necessary to construct a time-varying transverse broadening model using the cone-angle proxy extracted in Chapter~3.
Only after acquiring frame-by-frame matched reference data can the sensitivity to \(\alpha\) be reliably assessed.
